\documentclass[11pt]{article}
\usepackage{amsmath, amssymb}
\usepackage[margin=1in]{geometry}

\title{LaTeX in the Browser}
\author{Tectonic, compiled to WebAssembly}
\date{}

\begin{document}
\maketitle

\section{What this is}

The whole XeTeX engine is running in a Web Worker next to this editor.
Type in the pane on the left and the PDF on the right is rebuilt.

The interesting part is not that it compiles --- it is that it does not
recompile. The engine stays resident and keeps a ladder of checkpoints of
its own state. When you edit, it rewinds to the deepest checkpoint whose
inputs are still valid and resumes from there, so the preamble, the class
file and everything above your cursor are not touched.

\section{Mathematics}

Inline math like $e^{i\pi} + 1 = 0$ works, and so do displays:

\begin{equation}
  \int_{-\infty}^{\infty} e^{-x^{2}}\,dx = \sqrt{\pi}
\end{equation}

\begin{align}
  \nabla \times \mathbf{E} &= -\frac{\partial \mathbf{B}}{\partial t} \\
  \nabla \times \mathbf{B} &= \mu_0 \mathbf{J}
      + \mu_0 \varepsilon_0 \frac{\partial \mathbf{E}}{\partial t}
\end{align}

\section{Try it}

\begin{itemize}
  \item Edit this sentence and watch the latency counter above.
  \item Add a \verb|\section{}| and see the page reflow.
  \item Break something --- write \verb|\frac{1}| --- and read the log.
\end{itemize}

Edits near the end of the document are the cheapest, because there is more
above them to reuse. That is the number reported as \emph{reused} in the
toolbar.

\end{document}
